\section{Model background and governing equations}
\label{sec:model-background}

OpenIFS SURF represents the land surface as a set of independent vertical
columns coupled to the atmosphere through surface fluxes.  Each column is
partitioned into surface tiles, so the prognostic state is naturally written as
tile variables and layer variables,
\begin{equation}
  \mathbf{x} = \{\mathbf{T}_{\rm skin}, \mathbf{T}_{\rm soil},
  \mathbf{W}_{\rm soil}, S_{\rm snow}, f_{\rm snow}, W_{\rm skin},
  T_{\rm snow}, \rho_{\rm snow}, \alpha_{\rm snow}\}.
\label{eq:model-state}
\end{equation}
Here $\mathbf{T}_{\rm skin}$ is the tile skin temperature, $\mathbf{T}_{\rm
soil}$ and $\mathbf{W}_{\rm soil}$ are the four-layer soil temperature and
water stores, $S_{\rm snow}$ is snow water equivalent, $f_{\rm snow}$ is the
diagnosed snow-cover fraction, $W_{\rm skin}$ is intercepted canopy or skin
water, and $T_{\rm snow}$, $\rho_{\rm snow}$, and $\alpha_{\rm snow}$ are snow
temperature, density, and albedo.  Grid-box diagnostics are obtained by tile
area weighting,
\begin{equation}
  \bar{q} = \sum_{i=1}^{N_{\rm tile}} f_i q_i, \qquad
  \sum_{i=1}^{N_{\rm tile}} f_i = 1,
\label{eq:model-tilemean}
\end{equation}
where $f_i$ is the fractional area of tile $i$
\citep{ViterboBeljaars1995,Balsamo2009,Boussetta2021ECLand}.  This tiled structure is
preserved explicitly, and the snow and wet-skin tile fractions are updated
after each time step.

\subsection{Atmospheric forcing and offline surface coupling}

In coupled OpenIFS forecasts, SURF receives lower-atmosphere state variables,
near-surface exchange coefficients, and vertical diffusion boundary quantities
from the atmospheric physics.  In offline mode these quantities must
instead be constructed from meteorological forcing.  The forcing vector used by the
offline PyTorch implementation is
\begin{equation}
  \mathbf{y} = \{T_a, q_a, p_s, u_a, v_a, R_{\rm SW}^{\downarrow},
  R_{\rm LW}^{\downarrow}, P_{\rm rain}, P_{\rm snow}\},
\label{eq:model-forcing}
\end{equation}
with time interpolation or accumulation conventions matching the OpenIFS
offline configuration.  The surface stress and scalar increments are then closed by
the roughness-length update and the surface-layer vertical-diffusion
closure of the IFS physics package before the land update is invoked.
This closure is part of the offline PyTorch model only.
When SURF is coupled back into a complete atmospheric model, these quantities
are supplied by the atmospheric vertical-diffusion scheme rather than being
recomputed within SURF \citep{ECMWFIFSPartIV}.

\subsection{Surface energy balance}

For each active tile, the skin temperature is advanced from a surface energy
balance,
\begin{equation}
  C_i \frac{{\rm d}T_{s,i}}{{\rm d}t}
  = R_{n,i} - H_i - L_v E_i - G_i - M_i ,
\label{eq:model-seb}
\end{equation}
where $C_i$ is the effective surface heat capacity, $R_{n,i}$ is net radiation,
$H_i$ is sensible heat flux, $L_vE_i$ is latent heat flux, $G_i$ is the
conductive heat flux into the soil or snow pack, and $M_i$ is the latent heat
sink or source associated with melt and freezing.  The net radiation is
\begin{equation}
  R_{n,i} = (1-\alpha_i)R_{\rm SW}^{\downarrow}
  + \epsilon_i R_{\rm LW}^{\downarrow}
  - \epsilon_i \sigma T_{{\rm rad},i}^{4},
\label{eq:model-rn}
\end{equation}
where $\alpha_i$ is tile albedo, $\epsilon_i$ is emissivity, $\sigma$ is the
Stefan--Boltzmann constant, and $T_{{\rm rad},i}$ is the radiative skin
temperature used by the atmospheric model.  The original scheme solves the skin
temperature implicitly over the coupling time step, which is retained in the
PyTorch implementation \citep{ViterboBeljaars1995,ECMWFIFSPartIV}.

\subsection{Turbulent exchange}

The turbulent fluxes use a resistance or bulk-transfer form
\citep{Louis1979,ECMWFIFSPartIV},
\begin{equation}
  H_i = \rho_a c_p C_{H,i} |\mathbf{V}|(T_{s,i}-T_a),
  \qquad
  E_i = \rho_a C_{E,i} |\mathbf{V}|(q_{s,i}-q_a),
\label{eq:model-flux}
\end{equation}
where $\rho_a$ is air density, $c_p$ is the heat capacity of air, $C_{H,i}$ and
$C_{E,i}$ are stability- and roughness-dependent exchange coefficients,
$|\mathbf{V}|$ is the near-surface wind speed, and $q_{s,i}$ and $q_a$ are
surface and atmospheric specific humidity.  Momentum exchange is represented
analogously through the drag coefficient $C_D$.  In the tensor formulation
these algebraic coefficient evaluations are expressed as element-wise
differentiable maps, while the discrete regime switches for snow, bare soil,
vegetation, and wet-skin tiles are retained unchanged.

\subsection{Canopy interception and wet skin}

The intercepted water reservoir evolves according to
\begin{equation}
  \frac{{\rm d}W_{\rm skin}}{{\rm d}t}
  = P_{\rm int} - E_{\rm int} - D_{\rm int},
\label{eq:model-wskin}
\end{equation}
where $P_{\rm int}$ is intercepted precipitation, $E_{\rm int}$ is evaporation
from the wet canopy or skin reservoir, and $D_{\rm int}$ is drainage or excess
water released once the tile-dependent capacity is exceeded.  The reservoir is
bounded by a vegetation-dependent maximum storage $W_{\rm skin}^{\max}(LAI,
f_{\rm veg})$,
and the wet-skin fraction is diagnosed from $W_{\rm skin}/W_{\rm
skin}^{\max}$.  This diagnostic affects the tile fractions and therefore feeds
back onto the turbulent and radiative flux weighting
\citep{ViterboBeljaars1995,VanDenHurk2000}.

\subsection{Soil heat}

The soil thermal state follows a four-layer diffusion equation
\citep{ViterboBeljaars1995,ECMWFIFSPartIV},
\begin{equation}
  C_k \frac{\partial T_k}{\partial t}
  =
  \frac{\partial}{\partial z}
  \left(\lambda_k \frac{\partial T}{\partial z}\right) + S_k,
  \qquad k=1,\ldots,4 ,
\label{eq:model-soilheat}
\end{equation}
where $C_k$ is volumetric heat capacity, $\lambda_k$ is effective thermal
conductivity, and $S_k$ contains surface and phase-change source terms.  The
equation is discretized in finite-volume form with inter-layer conductive
fluxes $\lambda_{k+1/2}(T_{k+1}-T_k)/(z_{k+1}-z_k)$; the implicit tridiagonal
solve and the layer geometry are preserved from the reference implementation
\citep{ViterboBeljaars1995,ECMWFIFSPartIV}.

\subsection{Soil hydrology}

HTESSEL updates soil water through infiltration, evapotranspiration,
percolation, and runoff \citep{Balsamo2009}.  For layer $k$,
\begin{equation}
  \Delta z_k \frac{\partial \theta_k}{\partial t}
  =
  I_{k-1/2} - I_{k+1/2} - E_k - R_k ,
\label{eq:model-hydro}
\end{equation}
where $\theta_k$ is volumetric soil water, $I$ is vertical liquid-water flux,
$E_k$ is root-zone water extraction, and $R_k$ is surface or subsurface runoff.
The prognostic water mass $W_k = \rho_w \Delta z_k \theta_k$
is bounded by soil-type-dependent residual, wilting-point, field-capacity,
and saturation values.  These bounds keep the update numerically stable and
preserve the hydrological constraints of the reference model.

\subsection{Snow mass, snow temperature, density, and cover}

The snow pack is represented as a single-layer store with prognostic mass,
temperature, density, and albedo
\citep{ViterboBeljaars1995,Dutra2010Snow,ECMWFIFSPartIV}.  Snow water equivalent obeys
\begin{equation}
  \frac{{\rm d}S_{\rm snow}}{{\rm d}t}
  = P_{\rm snow} - E_{\rm snow} - M_{\rm snow},
\label{eq:model-swe}
\end{equation}
where $P_{\rm snow}$ is snowfall, $E_{\rm snow}$ is sublimation, and
$M_{\rm snow}$ is melt.  Snow temperature is advanced by the same surface
energy balance as Eq.~(\ref{eq:model-seb}), written on the snow tile with the
snow-pack heat capacity $C_{\rm snow}$ and the melt energy sink $L_f
M_{\rm snow}$.  Snow density evolves through compaction and fresh-snow updates, while snow
albedo is aged and refreshed by snowfall.  The active OpenIFS configuration
in this study uses the density-scaled snow-cover diagnostic,
\begin{equation}
  f_{\rm snow} =
  \min\left(1,\,10\,S_{\rm snow}/\rho_{\rm snow}\right),
\label{eq:model-fsnow}
\end{equation}
with very small snow fractions clipped to zero for tile-partition stability.
This relation is central to the interpretation of regional maps because full
snow cover can occur with moderate SWE when the density-scaled threshold is
exceeded \citep{ECMWFIFSPartIV,Boussetta2021ECLand}.

\subsection{Differentiable tensor formulation}
\label{sec:tensor-formulation}

The full time step is expressed as a single nonlinear operator
\begin{equation}
  \mathbf{x}^{n+1}
  =
  \mathcal{M}_{\Delta t}
  \left(\mathbf{x}^{n}, \mathbf{y}^{n:n+1}, \mathbf{s}; \boldsymbol{\theta}\right),
\label{eq:model-step}
\end{equation}
where $\mathbf{s}$ contains static land descriptors and $\boldsymbol{\theta}$
contains physical parameters or lookup-table values.  Most of the constituent
updates are element-wise differentiable maps.  The scheme, however, is not a
smooth dynamical system: rain--snow partitioning, tile activation, soil-moisture
bounds, wet-skin capacity overflow, freezing thresholds, and snow-cover
clipping are conditional statements that partition the
state space into regimes with distinct governing expressions.

The translation treats this structure as follows: every branch
is kept exactly as in the reference, and none is smoothed, relaxed, or
replaced by a differentiable surrogate.  The resulting operator is
differentiable almost everywhere, and reverse-mode differentiation through
it returns the exact derivative of the realized computation on the active
branch; at a regime crossing the derivative is the one-sided derivative of
that branch, a local linearization of the switched dynamics.  This choice
has a methodological consequence: the differentiated operator is the same
operator whose numerical equivalence is verified in
Sect.~\ref{sec:fortran-torch-equivalence}, so any gradient-based inference
refers to the discrete IFS physics itself rather than to a smoothed
reconstruction of it.  Smoothing the switches would yield better-behaved
gradients but would break the equivalence with the reference solution.  The
verification of Sect.~\ref{sec:gradient-calibration} shows that the
retained branches do not degrade the derivatives: multi-step reverse-mode
derivatives agree with centred finite differences to within
$4\times10^{-4}$ relative error, and the fraction of columns whose local
response is dominated by a regime transition is mapped and reported
(0.19--4.17\,\% depending on season) rather than omitted.
